\section{多源数据、预处理与统一符号}\label{sec:data}
\subsection{资料组成与研究用途}
小时资料覆盖MET、RADM、SOIL、FLUX四类，共24份年度CSV、63个测量字段和37个原始质量标记。统一时间轴包含52\,584个小时；数值记录的实际边界为2021年4月23日11时至2026年8月27日8时（原始UTC+8标签），各变量按其有效观测和质量轨使用。区域站元数据位置为$28.86^{\circ}$N、$85.29^{\circ}$E，海拔4140\,m，适合描述高寒山区的背景变化\cite{common-waternet}。背景排序采用2021—2025年对应日历日期前后30日，事件窗口保留24、72和168小时。

\begin{table}[htbp]
\centering\small
\caption{数据体系与分析支持。小时数和字段数表示记录结构，具体统计使用对应质量轨的有效样本。}\label{tab:common-data}
\begin{tabularx}{\linewidth}{p{2.05cm}p{2.65cm}Y}
\toprule
资料 & 结构与空间支持 & 主要用途\\
\midrule
MET & 22字段，21质量标记；1、2、4、10、20\,m梯度 & 温湿风状态、垂直差异、气压与降水探索\\
RADM & 4字段，4质量标记 & 上下行短波、长波与辐射背景\\
SOIL & 12字段，12质量标记；0.1—1.6\,m土层 & 含水率、土温、热通量及层间梯度\\
FLUX & 25字段，单列探索质量轨 & 湍流、风分量、声温及通量特征\\
GFS & 3个起报，f000—f096共579个选定字段 & 假设延迟和理想输入下的逐小时机制回放\\
遥感与地形 & SAR速度对、光学像元、COPDEM及对象几何 & 物源/几何约束、路径与传播地形\\
外部监测 & 13失稳组、1真实右删失站 & 变形加速、生存模型与时间区间分组评价\\
\bottomrule
\end{tabularx}
\end{table}

Copernicus GLO-30地形以可追溯DSM构建共同60和120\,m计算网格；源产品主体采集年代为2011—2015年，垂直参考EGM2008\cite{common-copdem}。遥感资料用于几何、表面状态和变化线索，地形用于条件传播计算。县域逐日温度资料独立保留长期背景角色。GFS按起报、预报时效和有效时刻组织，将气象预报输入与区域站历史观测分别管理\cite{common-gfs}。

\subsection{质量视图与时空预处理}
数值表、质量标记表和来源记录保留相同行序。记原始测量为$x_{v,t}$，数值有效掩膜为$m^{\mathrm{num}}_{v,t}$，原始质量标记为$q_{v,t}$，主轨定义为
\begin{equation}
 m^{(0)}_{v,t}=m^{\mathrm{num}}_{v,t}\,\mathbf 1(q_{v,t}=0),\qquad
 x^{(0)}_{v,t}=\begin{cases}x_{v,t},&m^{(0)}_{v,t}=1,\\ \mathrm{NA},&\text{其余位置}.\end{cases}
 \label{eq:common-qc}
\end{equation}
其中数值有效表示有限值并排除来源哨兵9999.9。原始文本缺测、哨兵、来源插补和质量类别分别保留。FLUX及降水采用原质量标记未随附的探索轨，通过纳入与排除比较评估其对因素排序的影响。此处理使质量选择能够复查，同时保留不同物理变量的符号与量纲。

小时来源按UTC+8约定换算为UTC，事件叙述明确标注显示时区；各窗口采用半开时间区间，并通过$\pm1$小时平移实际改变样本集合。县域逐日表保持其日期标签和空间聚合层级。模型计算统一使用UTC、米、秒或小时及明确的温度/压力单位，图表同时标明其物理支持。

空间分析使用WGS84/UTM 45N（EPSG:32645）。栅格地形、源面掩膜、路径与目标几何以版本关联；网格间比较固定总释放体积后调整厚度。对影像分别检查有效覆盖、分类掩膜和读入范围，对对象几何保留其产品或重建来源。物理阶段、资料用途与几何版本共同决定每项结果的解释范围。

\subsection{共享符号、信息集与评价单位}
\begin{table}[htbp]
\centering\small
\caption{跨问主要符号；分问专用参数在对应模型中定义。}\label{tab:common-symbols}
\begin{tabularx}{\linewidth}{p{2.0cm}Yp{2.6cm}}
\toprule
符号 & 含义 & 单位/支持\\
\midrule
$t,t_0,H$ & 时刻、预测发布时刻、预测时域 & UTC；$H=24,48,72$\,h\\
$t_{\mathrm m},t_{\mathrm r},\tau_j$ & 首次机械启动、整体释放、目标相对传播时延 & h或s，分别注明\\
$\boldsymbol\theta$ & 给定机制的结构、热、水及材料参数 & 各分量按物理量给出\\
$T,\alpha$ & 温度、热扩散率 & $^{\circ}$C、m$^2$/s\\
$\sigma_n,p$ & 法向应力、孔隙水压力 & kPa\\
$h,\boldsymbol u$ & 流深、深度平均水平速度 & m、m/s\\
$V,M$ & 体积、质量 & m$^3$、kg\\
$\mathcal I(t_0)$ & 发布时刻可使用的信息集 & 按观测/假设延迟分列\\
\bottomrule
\end{tabularx}
\end{table}

本文分别组织区域观测分析、外部站点评价、给定条件情景和数值基准。因素重采样以年份和过程为层级，时间预测以物理站点或失稳组为分组单位，传播以配对情景为比较单位。重复发布、传感器和目标响应均在原分组内汇总。后续章节使用这些单位报告样本数、区间覆盖和敏感性，使不同模型的比较具有一致基准。
\AIcredit
